\eventsection{INTRODU\c{C}\~AO}

Na Amaz\^onia, a rela\c{c}\~ao entre popula\c{c}\~oes locais e os ambientes aqu\'aticos \'e imediata. Lagos, rios e igarap\'es sustentam deslocamentos, abastecimento, pesca, atividades produtivas e rotinas comunit\'arias, especialmente em territ\'orios ribeirinhos. Por essa raz\~ao, altera\c{c}\~oes na qualidade da \'agua podem gerar impactos ambientais e sociais relevantes, atingindo tanto o equil\'ibrio ecol\'ogico quanto as condi\c{c}\~oes de vida das comunidades que dependem desses recursos.

Apesar dessa import\^ancia, o acompanhamento sistem\'atico da qualidade da \'agua ainda enfrenta barreiras pr\'aticas. Em muitas \'areas amaz\^onicas, as campanhas de coleta s\~ao espa\c{c}adas, o acesso aos pontos de medi\c{c}\~ao exige deslocamentos longos e a conectividade digital \'e inst\'avel ou inexistente. Nesse cen\'ario, solu\c{c}\~oes apoiadas apenas em transmiss\~ao cont\'inua de dados para servi\c{c}os remotos tendem a apresentar limita\c{c}\~oes operacionais e financeiras. A literatura aponta que sistemas baseados em Internet das Coisas e intelig\^encia artificial podem ampliar a capacidade de monitoramento, mas sua efic\'acia depende de adapta\c{c}\~ao \`as condi\c{c}\~oes reais de rede, log\'istica e uso em campo \cite{zulkifli2022,rahu2023,lowe2022}. Esse contexto favorece abordagens capazes de processar informa\c{c}\~oes no pr\'oprio local de uso. A intelig\^encia artificial de borda surge como alternativa promissora ao permitir infer\^encia local, menor depend\^encia de internet e maior autonomia dos dispositivos, caracter\'isticas valiosas para a realidade amaz\^onica \cite{aral2024}. \'E nesse horizonte que se insere este trabalho, voltado a solu\c{c}\~oes tecnol\'ogicas de baixo custo ajustadas \`a realidade amaz\^onica.

Neste artigo, apresenta-se uma proposta aplicada para o Lago do Serpa, em Itacoatiara, Amazonas, com foco no monitoramento offline de pH, oxig\^enio dissolvido e turbidez. A \^enfase n\~ao recai sobre o detalhamento interno da modelagem, mas sobre a demonstra\c{c}\~ao de uso de um sistema embarcado capaz de operar em plataformas acess\'iveis e em um cen\'ario territorial concreto.

Do ponto de vista do eixo tem\'atico do evento, a relev\^ancia da proposta est\'a em aproximar tecnologia e necessidade social. Em vez de tratar a qualidade da \'agua apenas como um problema de laborat\'orio ou de alto desempenho computacional, o estudo a enquadra como tema de vigil\^ancia ambiental, preven\c{c}\~ao de riscos e apoio \`a tomada de decis\~ao em regi\~oes remotas. Este trabalho deriva de um estudo anterior voltado \`a modelagem e \`a avalia\c{c}\~ao t\'ecnica de intelig\^encia artificial embarcada. Nesta vers\~ao, a prioridade recai sobre a aplica\c{c}\~ao pr\'atica, a leitura territorial dos resultados e o contexto socioambiental amaz\^onico. A principal contribui\c{c}\~ao deste trabalho est\'a na demonstra\c{c}\~ao de um sistema funcional, embarcado e offline, validado em m\'ultiplas plataformas e aplicado a um cen\'ario real da Amaz\^onia. Nesse sentido, a proposta se apresenta como alternativa para monitoramento ambiental de baixo custo, apoio \`a sa\'ude p\'ublica e suporte a comunidades ribeirinhas em regi\~oes remotas.

As principais contribui\c{c}\~oes deste trabalho podem ser sintetizadas em quatro pontos: o desenvolvimento de um sistema embarcado para monitoramento da qualidade da \'agua; a valida\c{c}\~ao da solu\c{c}\~ao em m\'ultiplas plataformas acess\'iveis; a demonstra\c{c}\~ao de aplicabilidade em um cen\'ario real da Amaz\^onia; e a proposi\c{c}\~ao de um uso voltado ao monitoramento ambiental de baixo custo em regi\~oes de baixa conectividade.
